\section{Exploratory Computing}
\label{sec:EC}

With the advent of Big Data, the increasing complexity and increase of easily available 
compute resources, the field of data analysis has been under an ever-increasing pressure 
to "make sense" of the abundant data.
This need has given rise to the methodology of \acf{EC}, which is an iterative 
process of data analysis which maintains a close to real-time dialogue between the
analyst and the data\footcite{diblasExploratoryComputingDraft2014}.

Unlike conventional data analysis methods which might follow a straighter path towards
an answer to a specific question, \ac{EC} aims to be a more open-ended process, where the eventual goal is
not just to find a numerical answer to a question but rather to gain a deeper understanding of the dataset.
Therefore, the tools in this field require a high degree of interactivity in order to
become a more natural extension of the analysts thought process\footcite{grangerJupyterThinkingStorytelling2021}.

\textbf{Computational Notebooks}

The defacto standard for scientific \ac{EC} is the \textit{Computational Notebook},
allowing for the direct integration of narrative text, code, and it's output like visualizations, tables and other data structures.
While in some domains other notebooks systems like \textit{Wolfram Mathematica} and \textit{Matlab} are used,
the most common notebook especially in the field of data science is the \textit{Jupyter Notebook}\footcite{perkelWhyJupyterData2018}.

\textit{Jupyter Notebooks} is a \ac{FOSS} server-(web)client application that exposes 
an interactive editor.
User can write code, documentation and visualizations into so-called cells. 
These cells are then executed on the server, which can be local or a remote compute resource,
and the output is displayed in the notebook.
As the computational environment is preserved across cells, 
with each cell being able to access the variables and functions defined in previous cells.
This lends itself to the \ac{REPL} style of programming, without the need to re-run the whole script\footcite{jupyter-docsProjectJupyterDocumentation}.
One of the key features of \textit{Jupyter Notebooks} is the ability to run code in multiple languages,
as the server on the  backend is build on a modular architecture that allows for the integration of \textit{kernels}.
These kernels are responsible for executing the code in the cells and returning the output to the notebook,
but are only limited by the \ac{API} spec of the notebook server\footcite{jupyter-docsArchitectureJupyterDocumentation}.
This has lead to a wide variety of kernels being available, with the most common ones being for \textit{Python}, \textit{R} and \textit{Julia},
but also supporting languages which are not traditionally associated with data analysis like \textit{Postsript}, \textit{PowerShell} and \textit{Brainfuck}\footcite{jupyter-docsJupyterKernels}. 

\pagebreak